\section{Balancing coefficients without adding variables}

The preceding arguments allow the fixed anchor integer
\[
 S_0=3^{2790}=c^6,\qquad c=3^{465}.
\]
The core produced by the selected source has 392 active rational
coordinates and 360 quadratic residuals. Its integer coefficients are
already explicit after specializing $S_0$. A diagonal rational change
of twelve existing coordinates makes those coefficients substantially
smaller without introducing a constant-chain witness.

\subsection{The exact change of coordinates}

Let $x_i$ denote a coordinate of the compiled source system and $y_i$
its replacement. The distinguished input is unchanged: $x_0=y_0=X$.
All unlisted coordinates are unchanged. For the listed coordinates set
$x_i=c^{e_i}y_i$ with the following exponents:
\begin{center}
\begin{tabular}{lll}
\toprule
Source coordinates & Number & $e_i$\\\midrule
The two coefficients of $h_S$ &2&1\\
The two stored coefficients of $x(S)$ &2&1\\
The two coefficients of $w=h_S\tau$ &2&2\\
The two coefficients of $f(S-P)$ &2&1\\
The three translated anchor witnesses $U,V,W$ &3&4\\
The materialized physical anchor $\Theta$ &1&4\\\bottomrule
\end{tabular}
\end{center}

The coordinate names and integer source identifiers are bound in the
machine-readable certificate. The coefficients of $h_S$ and $x(S)$
are intentionally scaled by the same factor, although the old identity
is $x(S)=h_S^2$. Their new defining equations incorporate that factor;
the new coordinates are decoded before being interpreted as the old
point functions. No semantic identity is silently assumed unchanged.

Since $c$ is a fixed nonzero rational integer, this diagonal map is a
bijection on $\mathbb Q^{392}$. Its inverse divides a listed coordinate
by $c^{e_i}$. It fixes the input and does not introduce any square-root
choice, new quantified parameter, or arithmetic hypothesis.

Let $R_j(x,S_0)$ be an integer quadratic residual of the compiled source.
Substitute $x_i=c^{e_i}y_i$ and $S_0=c^6$, combine like monomials, and
let $d_j\ne0$ be the signed integer content of the resulting row, with
the sign chosen to make the first coefficient positive. Set
\[
 Q_j(y)=\frac{R_j((c^{e_i}y_i)_i,c^6)}{d_j}.
\]
Every $Q_j$ is a primitive integer polynomial of degree at most two.
The delivered polynomial is exactly
\[
 P_{392}(y)=\sum_{j=1}^{360}Q_j(y)^2.
\]
Over $\mathbb Q$, a sum of squares vanishes if and only if each
summand vanishes. Thus
\[
 P_{392}(y)=0\quad\Longleftrightarrow\quad
 R_j((c^{e_i}y_i)_i,c^6)=0\quad\text{for every }j.
\]
Combining this equivalence with the diagonal bijection proves equality
of the existential projections to $X$. The new polynomial need not be
a single constant multiple of the old sum of squares: its individual
rows have been rescaled, which preserves their common zero set.

\subsection{Why the exponent drops}

Before the change, the fixed parameter occurs with exponent six in $c$.
After substitution, each affected row has a common power of $c$ whose
removal leaves a span at most two in its remaining $c$-exponents.
For example, the physical anchor row contains $\Theta$ and $N(w)$
at exponent four and $S_0N(B_0+X\sqrt2)$ at exponent six. Dividing
by its common factor leaves exponent two. In the translated anchor
row, the quadratic witness terms have exponent eight and the
$S_0U$ term has exponent ten, again leaving a span of two.

The shifted-point equations, the root chart equations, and the scalar
leaves have the same bounded span. All of these statements are checked
on the actual compiled rows. Consequently their squares involve at
most $c^4$ in the large coefficient, in place of the former $S_0^2=c^{12}$.

For the selected source, exact full expansion gives
\[
 \max|\operatorname{coeff}(P_{392})|
 =1+40532396646334464\,c^4
 =1+40532396646334464\,3^{1860}.
\]
This integer has 905 decimal digits and 3004 bits. The complete
quartic has 26,301 nonzero monomials, content one, and coefficient
one on $X^4$. Every coordinate occurs. The two constant-chain
coordinates used in the preceding coefficient-controlled construction
are therefore unnecessary at its former 978-digit cap.

The diagonal pattern was suggested by a linear optimization of
coefficient valuations. That floating-point search is not a verification
dependency. The selected pattern uses only the displayed integer
exponents; the delivered producer uses Python's standard library and
exact integers. A separately authored checker reconstructs all 360
substituted rows, verifies each nonzero row divisor and the input-fixed
diagonal map, independently expands the complete sum of squares,
and compares every coefficient of the text table. The search makes
no optimality claim.
